\documentclass[12pt]{article}
\usepackage{amsmath, amssymb}
\usepackage{geometry}
\geometry{margin=1in}

\title{CSE400 -- Fundamentals of Probability in Computing}
\author{Lecture L15--L18}
\date{}

\begin{document}

\maketitle

\section*{1. Why Study Joint Random Variables?}

\textbf{Concept}

In many real-world problems, we are interested in more than one random variable at the same time.

\textbf{Reasoning}

A single random variable captures only one aspect of a system. \\
Real systems depend on multiple interacting quantities. \\
Therefore, we must study their combined behavior $\rightarrow$ Joint Random Variables.

\subsection*{Example 1: Smart Transportation}

Let:
\begin{itemize}
\item $X =$ Vehicle Speed
\item $Y =$ Traffic Density
\end{itemize}

\textbf{Question}

Can speed be high when traffic density is high?

\textbf{Step-by-step reasoning (from slide intent)}

If traffic density increases $\rightarrow$ more vehicles on road \\
More vehicles $\rightarrow$ congestion increases \\
Congestion $\rightarrow$ reduces speed

Thus:
\begin{itemize}
\item Relationship exists between $X$ and $Y$
\item They are correlated random variables
\end{itemize}

\textbf{Applications}
\begin{itemize}
\item Adaptive traffic signals
\item Congestion prediction
\item Autonomous vehicle routing
\end{itemize}

\subsection*{Example 2: Wireless Network Performance}

Let:
\begin{itemize}
\item $X =$ Signal-to-Noise Ratio (SNR)
\item $Y =$ Data Throughput
\end{itemize}

\textbf{Question}

If SNR is high, will throughput always be high?

\textbf{Reasoning (as implied in slides)}

High SNR $\rightarrow$ better signal quality \\
But throughput also depends on:
\begin{itemize}
\item Network congestion
\item Interference
\end{itemize}

Hence:
High SNR does not guarantee high throughput

\textbf{Conclusion}

Variables are conditionally related \\
Need joint modeling

\textbf{Applications}
\begin{itemize}
\item 5G scheduling
\item Adaptive modulation
\item Network planning
\end{itemize}

\textbf{Other Examples}
\begin{itemize}
\item Weather: Rainfall, Temperature
\item Drone Delivery: Wind Speed, Delivery Time
\item Dice: Two random variables (requires joint PMF)
\end{itemize}

\section*{2. Discrete Case -- Joint PMF}

\textbf{Definition (implied from slides)}

Joint PMF gives probability of two discrete random variables taking specific values simultaneously.

\textbf{Reasoning}

For single RV $\rightarrow$ probability of one outcome \\
For two RVs $\rightarrow$ probability of pair of outcomes

Hence:
Need a function of two variables

\[
P(X = x, Y = y)
\]

\subsection*{Joint PMF Table Representation}

Rows $\rightarrow$ values of $X$ \\
Columns $\rightarrow$ values of $Y$

Each cell $\rightarrow P(X = x, Y = y)$

\textbf{Reasoning}

Since both variables vary independently in values, \\
Table helps visualize all combinations

\subsection*{Example: Rolling Two Dice}

Let:
\begin{itemize}
\item $X =$ outcome of first die
\item $Y =$ outcome of second die
\end{itemize}

\textbf{Step-by-step reasoning}

Each die has values: $\{1,2,3,4,5,6\}$ \\
Total outcomes: $6 \times 6 = 36$ \\
Each pair $(x,y)$ is equally likely:

\[
P(X = x, Y = y) = \frac{1}{36}
\]

\textbf{Example Questions}

\textbf{1. $X + Y = 7$}

Pairs:
\[
(1,6), (2,5), (3,4), (4,3), (5,2), (6,1)
\]

\[
P(X + Y = 7) = \frac{6}{36} = \frac{1}{6}
\]

\textbf{2. $X > Y$}

Count all pairs where first die $>$ second die

Total favorable outcomes = 15

\[
P(X > Y) = \frac{15}{36}
\]

\section*{3. Continuous Case -- Joint PDF}

\textbf{Definition (implied)}

Joint PDF describes probability over continuous variables:

\[
f_{X,Y}(x,y)
\]

\textbf{Key Idea}

Probability is not at a point \\
It is over a region

\subsection*{Geometric Interpretation}

\[
P((X,Y) \in R) = \iint_R f_{X,Y}(x,y)\,dx\,dy
\]

\textbf{Reasoning}

Continuous case $\rightarrow$ infinite possible values

So:
\begin{itemize}
\item Probability at a point = 0
\item Use area (integration) over region
\end{itemize}

\subsection*{Worked Example (structure)}

\begin{enumerate}
\item Identify region of support
\item Use limits of integration
\item Integrate joint PDF over region
\item Compute required probability
\end{enumerate}

\section*{4. Joint Cumulative Distribution Function (Joint CDF)}

\textbf{Definition}

\[
F_{X,Y}(x,y) = P(X \le x, Y \le y)
\]

\textbf{Reasoning}

CDF accumulates probability \\
For joint case: accumulate over 2D region

\section*{5. Joint CDF -- Continuous Case}

\[
F_{X,Y}(x,y) = \int_{-\infty}^{x} \int_{-\infty}^{y} f_{X,Y}(u,v)\,du\,dv
\]

\textbf{Step-by-step reasoning}

\begin{enumerate}
\item Fix upper bounds $x, y$
\item Consider all values less than them
\item Integrate PDF over that region
\end{enumerate}

\section*{6. Joint CDF -- Discrete Case}

\[
F_{X,Y}(x,y) = \sum_{u \le x} \sum_{v \le y} P(X=u, Y=v)
\]

\textbf{Reasoning}

Replace integration with summation \\
Because variables are discrete

\section*{7. Properties of Joint CDF}

\begin{enumerate}
\item \textbf{Non-decreasing}

If $x_1 \le x_2$, $y_1 \le y_2$, then:
\[
F(x_1, y_1) \le F(x_2, y_2)
\]

\textbf{Reasoning}: Larger region $\rightarrow$ more accumulated probability

\item \textbf{Limits}
\[
\lim_{x,y \to \infty} F(x,y) = 1
\]

\textbf{Reasoning}: Entire probability space covered

\item \textbf{Non-negative}
\[
F(x,y) \ge 0
\]

\item \textbf{Boundary behavior}

At very small values $\rightarrow 0$ \\
At very large values $\rightarrow 1$
\end{enumerate}

\section*{8. Comprehensive Understanding}

\textbf{Flow of Concepts}

\begin{enumerate}
\item Real-world motivation $\rightarrow$ multiple RVs
\item Discrete case $\rightarrow$ Joint PMF
\item Continuous case $\rightarrow$ Joint PDF
\item Accumulation $\rightarrow$ Joint CDF
\item Properties $\rightarrow$ behavior of distributions
\end{enumerate}

\textbf{Core Insight}

Systems depend on relationships between variables

Joint distributions capture:
\begin{itemize}
\item Dependence
\item Correlation
\item Combined probabilities
\end{itemize}

\section*{End of Lecture}

Thank You!

\end{document}
